\scanpage{62}
\mspage{\handwritten{- 34 -}}
La fonction égale à $\handwritten{f}$ dans $A$, à $\handwritten{g}$ dans $B$, est alors holomorphe dans $A\handwritten{\cup}B$ ; on dit que c'est un \uline{prolongement analytique} de $\handwritten{f}$ (ou de $\handwritten{g}$) à $A\handwritten{\cup}B$.

Le th.1 entraîne le résultat plus général suivant :

PROPOSITION 1.- \uline{Soit} $(A_{\handwritten{\alpha}})$ \uline{une famille de domaines dans} $K^p$, \uline{telsque}, \uline{pour tout couple d'indices} $\handwritten{\alpha,\beta}$, $A_{\handwritten{\alpha}}\handwritten{\cap}A_{\handwritten{\beta}}$ \uline{soit vide ou soit un domaine. Pour tout indice} $\handwritten{\alpha}$, \uline{soit} $\handwritten{f_\alpha}$ \uline{une fonction à valeurs dans} $E$, \uline{définie et holomorphe dans} $A_{\handwritten{\alpha}}$ ; \uline{on suppose que pour tout couple d'indices} $\handwritten{\alpha,\beta}$ \uline{tel que} $A_{\handwritten{\alpha}}\handwritten{\cap}A_{\handwritten{\beta}}\handwritten{\neq\varnothing}$, \uline{il existe un ensemble ouvert non vide contenu dans} $A_{\handwritten{\alpha}}\handwritten{\cap}A_{\handwritten{\beta}}$ \uline{et dans lequel} $\handwritten{f_\alpha(x)=f_\beta(x)}$. \uline{Dans ces conditions, il existe une.fonction} $\handwritten{f}$, \uline{holomorphe dans} $A=\handwritten{\bigcup_\alpha}A_{\handwritten{\alpha}}$, \uline{et telle que} $\handwritten{f(x)=f_\alpha(x)}$ \uline{pour tout} $\handwritten{x\in}A_{\handwritten{\alpha}}$ \uline{et tout} $\handwritten{\alpha}$.

En effet, pour tout $\handwritten{x\in}A$, les valeurs de $\handwritten{f_\alpha(x)}$ sont les mêmes pour tout $\handwritten{\alpha}$ tel que $\handwritten{x\in}A_{\handwritten{\alpha}}$, en vertu du th.1.

\uline{\handwritten{2}}.\quad \uline{Domaines d'existence ma\handwritten{-}ximaux}.\quad \uline{Points singuliers}.

Soit $A$ un domaine dans $K^p$, et $\handwritten{f}$ une fonction holomorphe dans $A$. Considérons l'ensemble $\handwritten{\mathscr{M}}$ des fonctions holomorphes définies dans un domaine contenant $A$, et \uline{prolongeant} $\handwritten{f}$ ; il est immédiat que $\handwritten{\mathscr{M}}$, ordonné par la relation ``$\handwritten{g}$ est une restriction de $\handwritten{h}$'' est \uline{inductif} ; en vertu du th. de Zorn, il admet donc au moins un élément maximal $\handwritten{f_{\unsure{0|o}}}$ ; \margin{\unsure{i}}\unsure{.} on dit que $\handwritten{f_{0}}$ est un \uline{prolongement analytique maximal} de $\handwritten{f}$, et que le domaine $A_{\handwritten{0}}$ où est définie $\handwritten{f_{0}}$ est un \uline{domaine d'existence maximal} pour un prolongement analytique de $\handwritten{f}$ (ou par abus de langage, pour $\handwritten{f}$ elle-même).

Si $A_{\handwritten{0}}$ est un domaine d'existence maximal pour $\handwritten{f}$, il résulte aussitôt du th.1 que le prolongement maximal de $\handwritten{f}$ dans $A_{\handwritten{0}}$ est \uline{unique}. Si $A_{\handwritten{0}}=K^p$ il n'y a qu'un seul domaine d'existence maximal pour $\handwritten{f}$,
% archive sheet number: 61
